\chapter{Conclusioni}
\label{cap:5}

Questo lavoro ha confrontato due strategie di fine-tuning per il
trasferimento di capacità di ragionamento a modelli linguistici di
piccola taglia: il Reflection Tuning, basato sull'addestramento su
sequenze con ragionamento esplicito, e la Knowledge Distillation
response-based, basata sull'imitazione delle risposte di un modello
teacher più capace.

A partire dal dataset pubblico OpenThoughts-114k, sono stati estratti
700 campioni di training e 80 di valutazione. Il dataset di
distillation è stato generato interrogando Gemini Flash come teacher
su ciascuna delle 700 domande. I due dataset sono stati usati per
addestrare separatamente due versioni di Llama~3.2~1B-Instruct tramite
QLoRA, mantenendo identici tutti gli iperparametri. La valutazione
finale su 20 domande di ragionamento matematico è stata condotta con
Gemini Flash secondo il paradigma LLM-as-a-Judge \cite{zheng2023judging}.

I risultati mostrano una netta prevalenza del modello addestrato con
knowledge distillation: su 19 domande valide, il modello di
distillation vince in 14 casi (73.7\%), con punteggi medi superiori
in tutte e tre le metriche (Correctness 3.89 vs.\ 2.74, Reasoning
Quality 3.74 vs.\ 3.05, Clarity 4.11 vs.\ 2.42). Il vantaggio è
particolarmente marcato sulla Clarity ($\Delta = -1.68$), mentre si
riduce sulla Reasoning Quality ($\Delta = -0.68$), l'unica dimensione
in cui il modello reflection mostra risultati competitivi: nei cinque
casi in cui prevale, il blocco \texttt{<thinking>} sembra fungere da
meccanismo di auto-correzione su problemi multi-passo. Questi risultati
sono coerenti con la letteratura: a 1B parametri, l'imitazione diretta
di un teacher capace produce un comportamento più stabile rispetto
all'apprendimento di strutture di ragionamento complesse, che richiedono
scala e dati significativamente maggiori \cite{taori2023alpaca,wei2022chain}.

I risultati devono essere letti considerando i vincoli pratici
dell'intero lavoro. L'addestramento su Google Colab con GPU Tesla T4
gratuita ha imposto un tetto alla dimensione del modello utilizzabile:
architetture di scala superiore, su cui i vantaggi del reflection
tuning sono documentati in letteratura \cite{deepseekr1_2025}, erano
fuori portata. Il dataset di 700 campioni è stato vincolato dalla quota
gratuita dell'API Gemini (15 richieste/minuto) e dai limiti di sessione
di Colab; la valutazione su 20 domande consente di osservare tendenze
qualitative ma non conclusioni statisticamente significative. Va inoltre
precisato che la knowledge distillation implementata è di tipo
response-based --- lo student impara a imitare le risposte finali ---
e non incorpora i soft targets della formulazione classica
\cite{hinton2015distilling}, inaccessibili tramite API commerciali.
Il confronto è quindi più precisamente descritto come \emph{SFT su
dati con ragionamento esplicito} vs.\ \emph{SFT su risposte di un
modello capace}. Infine, il giudice automatico (Gemini Flash) coincide
con il teacher del modello di distillation, introducendo un potenziale
bias stilistico non quantificabile in assenza di valutatori indipendenti.

Per approfondimenti futuri, replicare l'esperimento con modelli di
scala maggiore (7B parametri o superiori) permetterebbe di verificare
se il vantaggio del reflection tuning emerge in modo più consistente.
Ampliare il dataset e il test set --- includendo benchmark standardizzati
come GSM8K \cite{cobbe2021gsm8k} o MATH \cite{hendrycks2021math} ---
consentirebbe una valutazione più robusta e confrontabile con la
letteratura. Sarebbe inoltre interessante esplorare un approccio ibrido
a due stadi: prima il teacher per acquisire correttezza di base, poi
dati con ragionamento esplicito per sviluppare capacità di auto-correzione.

Nonostante i vincoli, questo lavoro dimostra la fattibilità di una
pipeline completa di fine-tuning comparativo su hardware consumer-grade.
Il risultato principale --- la superiorità della distillation
response-based nelle condizioni adottate --- non confuta il reflection
tuning, ma indica che scala del modello e dimensione del dataset sono
fattori determinanti per sfruttare efficacemente il ragionamento
esplicito nel training. In altre parole, insegnare a un modello piccolo
a \emph{pensare ad alta voce} richiede che quel modello abbia già la
capacità di pensare: a 1B parametri, con 700 esempi, quella capacità
non è ancora sufficiente a tradursi in un vantaggio misurabile.
